\documentclass[10pt,a4paper]{amsart}
\usepackage[margin=19mm]{geometry}
\usepackage{amsmath,amssymb,mathtools}
\usepackage[scr=rsfs,cal=euler]{mathalfa}
\usepackage{xfrac}
\usepackage[unicode,hidelinks,bookmarks=false]{hyperref}
\newcommand{\ie}{i.e.\ }
\newcommand{\cf}{cf.\ }
\newcommand{\resp}{respectively\ }
\newcommand{\Subsection}{\subsection}
\providecommand{\nobreakdash}{\nobreak\hbox{-}}
\setlength{\emergencystretch}{2em}
\multlinegap0pt
\usepackage{amssymb}
\newtheorem{corollary}{Corollary}
\newtheorem{lemma}{Lemma}
\newtheorem{theorem}{Theorem}
\newcommand{\Dbb}{\mathbb D}
\newcommand{\Tbb}{\mathbb T}
\newcommand{\al}{\alpha}
\newcommand{\aream}{A}
\newcommand{\bneva}{\mathcal{N}}
\newcommand{\clk}{\sigma}
\newcommand{\cph}{C_\varphi}
\newcommand{\diin}{\Theta}
\newcommand{\dn}{\mathbb{D}^n}
\newcommand{\er}{\varepsilon}
\newcommand{\kknl}{\mathbf{k}}
\newcommand{\knl}{k}
\newcommand{\meal}{m_n}
\newcommand{\neva}{N}
\newcommand{\ph}{\varphi}
\newcommand{\spr}{\Sigma}
\newcommand{\tn}{\mathbb{T}^n}
\newcommand{\za}{\zeta}
\title{Composition operators between model and Hardy spaces}
\author{Evgueni Doubtsov}
\date{Source: arXiv:2604.04054v1 (CC BY 4.0)}
\begin{document}
\maketitle

\noindent\emph{Source and licence.} Composition operators between model and Hardy spaces. Version arXiv:2604.04054v1 (CC BY 4.0), distributed by arXiv under the Creative Commons Attribution 4.0 International licence, http://creativecommons.org/licenses/by/4.0/, as recorded in the arXiv licence metadata for this exact version and shown on the abstract page https://arxiv.org/abs/2604.04054v1. Full licence notice: \texttt{licenses/arxiv-2604.04054-CC-BY-4.0.txt}.\par\smallskip

\noindent\textit{Editorial scope.} Counted unit: the article's theorem T\_2 (source0.tex lines 421--428). The article's complete original proof of it is delivered with it (source0.tex lines 429--537, one \texttt{proof} environment of 109 lines). No mathematics is written, repaired or re-derived: the statement and the proof are the article's own lines, in the article's own notation, and no retained line is substituted or re-flowed.  Retained uncounted as the source-local context the proof needs: lines 194--197; lines 216--224; lines 241--248; lines 243--245; lines 254--266; lines 260--263; lines 377--384.  Nothing was omitted from any retained span, so the package carries no omission record.  The retained lines cite 3 works, whose entries are transcribed verbatim from the article's own bibliography source in the e-print and delivered with short editorial labels recorded in the item's reference mapping.  No retained line contains a figure environment, an image-inclusion command or drawing code, so no image is delivered and the package declares no asset dependency.  Editorial formatting only, in addition: an AMS \texttt{amsart} preamble that restates the article's own declarations and macros used by the retained lines, namely the environments \texttt{corollary}, \texttt{lemma}, \texttt{theorem} on separate counters, together with the standard packages those lines need.  The retained environments print in this document as Corollary 1, Lemma 1, Theorem 1 in order of appearance, whereas the article prints the same items with the numbering of its own full document.  The complete native source of the article as distributed in the arXiv e-print for this exact version is bundled unchanged as \texttt{sources/197858/source0.tex}.
\begin{equation}\label{e_Stntn_d}
\|f\circ \ph\|^2_{H^2(\dn)} = |f(\ph(0))|^2 + 2\int_{\Dbb} |f^\prime(w)|^2
\left(\int_{\tn} \neva_{\ph_\za}(w)\,  d\meal(\za) \right)\, d\aream(w).
\end{equation}

\begin{corollary}\label{c_subharm}
Let $w\in\Dbb$ and $\ph: \dn\to\Dbb$ be a holomorphic function.
Let $\Delta$ be a disk centered at $w$ and such that $\ph(0)\notin \Delta$.
Then
\begin{equation}\label{e_subharm}
\int_{\tn} \neva_{\ph_\za}(w)\,d\meal(\za) \le \frac{1}{\aream(\Delta)}\int_{\Delta}
\left(\int_{\tn} \neva_{\ph_\za}(z)\,  d\meal(\za) \right)\, d\aream(z).
\end{equation}
\end{corollary}

\begin{lemma}[{\cite[Lemma~1]{LM13}}]\label{l_LM}
Let $\{w_q\}_{q=1}^\infty \subset\Dbb$ be a sequence such that $|w_q|\to 1$ as $q\to \infty$ and
\begin{equation}\label{e_LM}
|\diin (w_q)| < a
\end{equation}
for some parameter $a\in (0,1)$. Then
$\knl_{w_q}/\|\knl_{w_q}\| \overset{w^\ast}\longrightarrow 0$ as $q\to \infty$.
\end{lemma}

\begin{equation}
|\diin (w_q)| < a
\end{equation}

\begin{theorem}\label{t_cph_tri}
Let $\ph: \dn\to\Dbb$, $n\ge 1$, be a holomorphic function and $\diin: \Dbb \to \Dbb$ be an inner function
such that $K_\diin$ is infinite dimensional.
Then the following properties are equivalent.
\begin{itemize}
  \item [(i)] One has
  \begin{equation}\label{e_nevaTo0}
%\bneva(w) :=
\int_{\tn}\neva_{\ph_\za}(w) \frac{1-|\diin(w)|}{1-|w|}\, d\meal(\za)\to 0 \quad \textrm{as}\ |w|\to 1-.
\end{equation}
  \item [(ii)] The operator $\cph: K_\diin \to H^2(\dn)$ is compact.
\end{itemize}
\end{theorem}

  \begin{equation}
%\bneva(w) :=
\int_{\tn}\neva_{\ph_\za}(w) \frac{1-|\diin(w)|}{1-|w|}\, d\meal(\za)\to 0 \quad \textrm{as}\ |w|\to 1-.
\end{equation}

\begin{lemma}[{\cite[Section~5]{VT88}}]\label{l_A}
Let $\diin$ be a one-component inner function and $\al\in\Tbb$. Then the following properties are equivalent.
\begin{enumerate}
  \item[$\bullet$] $\al\in\spr(\diin)$;
  \item[$\bullet$] $\liminf_{w\to\al} |\diin(w)| <1$;
  \item[$\bullet$] $\liminf_{r\to 1-} |\diin(r\al)| <1$.
\end{enumerate}
\end{lemma}

\begin{theorem}\label{T_2}
Let $\diin$ be a one-component inner function such that $K_\diin$ is infinite dimensional
and let $\ph: \dn\to \Dbb$ be a holomorphic function, $n\ge 1$.
Then $\cph: K_\diin \to H^2(\dn)$ is a compact operator if and only if
\begin{equation}\label{e_clark}
\|\clk_\al^s\| =0 \quad \textrm{for all $\al\in\spr(\diin)$}.
\end{equation}
\end{theorem}

\begin{proof}
Assume that $\cph$ is a compact operator.
Consider a point $\al\in\spr(\diin)$.
By Lemma~\ref{l_A}, there exists a sequence $\{r_q\}$ such that $r_q\nearrow 1-$ and
\begin{equation}\label{e_less1}
  \lim_{q\to \infty} |\diin(r_q\al)| <1.
\end{equation}

Lemma~\ref{l_LM} guarantees that
\[
\int_{\tn}\frac{|1-\diin(\ph(\za))\overline{\diin}(r_q\al)|^2}{|1- r_q\ph(\za)\overline{\al}|^2} \frac{1-|r_q|^2}{1-|\diin(r_q\al)|^2}\,d\meal(\za)
= \frac{\|\cph \knl_{r_q\al}\|^2_{2}}{\|\knl_{r_q\al}\|^2_2} \to 0\quad \textrm{as}\ q\to\infty.
\]
By \eqref{e_less1}, we obtain
\begin{equation}\label{e_null}
% \|\cph \mathbf{k}_{r_q\al}\|^2_{2} =
  \int_{\tn}\frac{1-|r_q|^2}{|1- r_q\ph(\za)\overline{\al}|^2}
\to 0\quad \textrm{as}\ q\to\infty.
\end{equation}
Now, we argue as in \cite{CM97} for $n=1$.
One has
\[
\frac{1-|r_q|^2}{|1- r_q\ph(\za)\overline{\al}|^2}
= \frac{1-|r_q \ph(\za)|^2}{|1- r_q\ph(\za)\overline{\al}|^2}
-  |r_q|^2 \frac{1-|\ph(\za)|^2}{|1- r_q\ph(\za)\overline{\al}|^2}.
\]
Observe that the pluriharmonic function
\[
\frac{1-|r_q \ph(\za)|^2}{|1- r_q\ph(\za)\overline{\al}|^2}
\]
is bounded on the polydisk $\dn$,
thus,
\[
\int_{\tn}\frac{1-|r_q \ph(\za)|^2}{|1- r_q\ph(\za)\overline{\al}|^2}\, d\meal(\za) =
\frac{1-|r_q \ph(0)|^2}{|1- r_q\ph(0)\overline{\al}|^2}\overset{q\to \infty}\longrightarrow  \|\clk_\al\|.
\]
On the other hand, by the monotone convergence theorem,
\[
\int_{\tn}|r_q|^2 \frac{1-|\ph(\za)|^2}{|1- \ph(\za) r_q\overline{\al}|^2}\, d\meal(\za) \overset{q\to \infty}\longrightarrow
\int_{\tn}\frac{1-|\ph(\za)|^2}{|1- \ph(\za)\overline{\al}|^2}\, d\meal(\za) = \|\clk_\al^a\|.
\]
Applying \eqref{e_null}, we obtain
\[
\|\clk_\al^s\| = \|\clk_\al\| - \|\clk_\al^a\| =0,
\]
as required.

To prove the reverse implication, assume that condition \eqref{e_clark} is satisfied.
By Theo\-rem~\ref{t_cph_tri}, it suffices to verify property \eqref{e_nevaTo0}.
Suppose that \eqref{e_nevaTo0} does not hold.
Then there exist a sequence $\{w_q\}_{q=1}^\infty$ and a point $\al\in\Tbb$ such that
$w_q\to \al$ as $q\to \infty$ and
\[
\int_{\tn}\neva_{\ph_\za}(w_q) \frac{1-|\diin(w_q)|}{1-|w_q|}\, d\meal(\za) > c_0 > 0.
\]
By Littlewood's inequality,
\[
\neva_{\ph_\za}(w) \le C (1-|w|), \quad\za\in\tn,\ \frac{1+|\ph(0)|}{2} < |w| <1.
\]
Therefore, property \eqref{e_LM} holds. Also, $\al\in\spr(\diin)$ by Lemma~\ref{l_A} and
\begin{equation}\label{e_NgC}
\int_{\tn} \frac{\neva_{\ph_\za}(w_q)}{1-|w_q|}\, d\meal(\za) > c > 0
\end{equation}
for all sufficiently large $q$.


Now, we consider the bounded operator $\cph: H^2(\Dbb) \to H^2(\dn)$.
Recall that the reproducing kernel for $H^2(\Dbb)$ is defined by
\[
\kknl_w (z) = \frac{1}{1- z\overline{w}}, \quad \|\kknl_w\|_2^2 = \frac{1}{1-|w|^2}.
\]
Let $D_\er(w) = \{z\in\Dbb: |z-w| < \er|1- z\overline{w}|\}$, that is,
let $D_\er(w)$ denote the pseudohyperbolic $\er$-disk centered at  $w\in \Dbb$.
Observe that
\begin{equation}\label{e_kknl}
|\kknl_w^\prime (z)| \ge C_\er \frac{|w_q|}{(1-|w_q|)^2}, \quad z\in D_\er(w).
\end{equation}
Sequentially applying equality \eqref{e_Stntn_d}, estimate~\eqref{e_kknl} and Corollary~\ref{c_subharm}, we obtain
\begin{equation}\label{e_rker_estim}
\begin{aligned}
\frac{\|C_\ph \kknl_{w_q}(z)\|^2}{\|\kknl_{w_q}\|^2}
    &\ge \frac{1}{\|\kknl_{w_q}\|^2} \int_{\Dbb} |\kknl^\prime_{w_q}(z)|^2
\left(\int_{\tn} \neva_{\ph_\za}(z)\,  d\meal(\za) \right)\, d\aream(z)\\
    &\ge \int_{D_\er(w_q)} \frac{C}{(1- |w_q|^2)^{3}}
\left(\int_{\tn} \neva_{\ph_\za}(z)\,  d\meal(\za) \right)\, d\aream(z)\\
    &\ge \frac{C_\er}{1- |w_q|^2} \int_{\tn} \neva_{\ph_\za}(w_q)\,d\meal(\za).
\end{aligned}
\end{equation}


On the one hand, using \eqref{e_rker_estim} and \eqref{e_NgC}, we conclude that
\[
\frac{\|\cph \kknl_{w_q}\|_2^2}{\|\kknl_{w_q}\|_2^2}
\ge C \int_{\tn} \frac{\neva_{\ph_\za}(w_q)}{1-|w_q|}\, d\meal(\za) > C\cdot c > 0.
\]
On the other hand, applying Fatou's lemma and condition \eqref{e_clark}, we obtain
\[
\begin{split}
   &\limsup_{q\to \infty} \frac{\|\cph \kknl_{w_q}\|_2^2}{\|\kknl_{w_q}\|_2^2}
   \\
&\le \limsup_{q\to \infty} \int_{\tn} \frac{1-|w_q\ph(\za)|^2}{|1- \overline{w_q}\ph(\za)|^2} \, d\meal(\za)
- \liminf_{q\to \infty} \int_{\tn} |w_q|^2 \frac{1-|\ph(\za)|^2}{|1- \overline{w_q}\ph(\za) |^2}\, d\meal(\za) \\
&\le \frac{1-|\ph(0)|^2}{|1- \overline{\al}\ph(0)|^2}
- \int_{\tn} \frac{1-|\ph(\za)|^2}{|1- \overline{\al}\ph(\za) |^2}\, d\meal(\za) \\
     &= \|\clk_\al\| - \|\clk_\al^a\| = \|\clk_\al^s\| =0,
\end{split}
\]
a contradiction. Therefore, the required property \eqref{e_nevaTo0} holds.
\end{proof}

\begin{thebibliography}{99}
\bibitem[CM97]{CM97} J.~A. Cima and A.~L. Matheson, \emph{Essential norms of composition operators and {A}leksandrov measures}, Pacific J. Math. \textbf{179} (1997), no.~1, 59--64. \MR{1452525}
\bibitem[LM13]{LM13} Yu.~I. Lyubarskii and E.~Malinnikova, \emph{Composition operators on model spaces}, Recent trends in analysis, Theta Ser. Adv. Math., vol.~16, Theta, Bucharest, 2013, pp.~149--157. \MR{3411049}
\bibitem[VT88]{VT88} A.~L. Vol'berg and S.~R. Treil', \emph{Imbedding theorems for the invariant subspaces of the backward shift operator}, J. Soviet Math. \textbf{42} (1988), no.~2, 1562--1572. \MR{849293}
\end{thebibliography}
\end{document}
